\documentclass[12p, oneside]{book}
\usepackage[T1]{fontenc}
\usepackage[polish]{babel}
\usepackage{color}
\usepackage[table, usenames, dvipsnames]{xcolor}
\usepackage{hyperref}
\hypersetup{
    colorlinks=true, 
    linktoc=all,    
    linkcolor=blue,
}



\newcommand{\steno}[1]{\textbf{\textsf{\textcolor{violet}{ #1}}}}

\newcommand{\stenoExc}[1]{\textbf{\textsf{\textcolor{red}{ #1}}}}

\newenvironment{tabela2}
	{
		\rowcolors{1}{cyan}{Emerald}
		\begin{tabular}{c c}
	}
	{
		\end{tabular}
	}


\setcounter{secnumdepth}{0}

\begin{document}

\frontmatter

\title{Zachodniosłowiańska Stenografia Maszynowa\\ 
\normalsize Kurs natychmiastowego tworzenia tekstu z wykorzystaniem stenografi maszynowej}

\author{Priscilla Trillo}
\maketitle
\newpage


\tableofcontents

\chapter{Wstęp}

\pagebreak
\section{Od autora}
Witaj. Mam na imię Priscilla Trillo i jestem odpowiedzialna za powstanie książki, którą właśnie czytasz. Jestem profesjonalnym stenografem ze Stanów Zjednoczonych oraz tak się składa, że posiadam również pewne doświadczenie w zakresie języków słowiańskich. Ponad 40 lat temu zaczęłam się ich uczyć na własną rękę, a następnie kontynuowałam tę naukę na uniwersytecie. Jestem związana ze stenografią od 1996 roku; uczyłam stenografii maszynowej w języku angielskim na lokalnej uczelni. Jako stenograf pracowałam w Tanzanii dla Organizacji Narodów Zjednoczonych w Międzynarodowym Trybunale Karnym dla Rwandy (MTKR).

Chciałabym wyrazić wdzięczność Krzysztofowi Smirnowowi za wspaniałe grafiki zachodniosłowiańskiej klawiatury stenograficznej, jak również za zbieranie informacji na temat kombinacji spółgłosek występujących w języku polskim.

Wszystkie firmy, produkty i znaki towarowe wzmiankowane w tej książce są własnością odpowiednich właścicieli. Nie mam żadnego interesu finansowego powiązanego z którymkolwiek z tych produktów. Ich umieszczenie w książce ma na celu ułatwienie używania opisywanych metod. Czytelnicy tej książki mogą używać dowolnych lub wszystkich wspomnianych produktów zgodnie ze swoją wolą.

Książka została napisana po angielsku i celowo nie została przetłumaczona na żaden z omawianych w niej języków, ponieważ w razie potrzeby pomoc na Discordzie Plovera będzie dostępna właśnie po angielsku, gdyż tym językiem włada większość jego użytkowników. W przypadku zadania pytania w innym języku, może się zdarzyć, że konieczne będzie poczekać na odpowiedź przez dłuższy czas. \footnote{Jak widać, książka ta została przetłumaczona. Wszystkie uwagi dotyczące tłumaczenia oraz zmian względem oryginału znajdują się w sekcji Od Tłumacza}

W przypadku potrzeby kontaktu ze mną, można swobodnie do mnie napisać na adres trillopriscilla@gmail.com. Jestem członkiem Discorda Plovera; jednak jeżeli chcesz się ze mną tam skontaktować, to musisz mnie otagować, @Plover-Trillo. Nie sprawdzam go regularnie, ale zostanę powiadomiona, gdy zostanę otagowana w poście.


Z wyrazami szacunku,


Priscilla Trillo

\newpage
\section{Od tłumacza}

\newpage
\section{Przedmowa}
Ten podręcznik nie ma na celu nauczenia Czytelnika żadnego z opisywanych języków. Wręcz przeciwnie, zakłada się znajomość jednego lub wszystkich z omawianych języków. Nie będą omawiane zasady gramatyczne. Celem tej książki jest nauczenie szybkiego kodowania tekstu w celu natychmiastowego wyświetlenia lub tłumaczenia symultanicznego. Metoda zaprezentowana w tej książce pozwoli osiągnąć oba te cele. Na pierwszy rzut oka zagadnienie wdaje się być dość proste, ale osiągnięcie biegłości wymaga czasu i poświęceń.

Zdecydowałam się na umieszczenie wszystkich języków zachodniosłowiańskich razem w jednej książce, ponieważ dwa z nich są wzajemnie zrozumiałe, a trzeci nie jest od pozostałych dwóch tak bardzo odmienny. W rzeczywistości, zobaczenie jak jeden język radzi sobie z pewnym typem słów może uczynić łatwiejszym zrozumienie, dlaczego inny język, choć odrobinę inny, może być kodowany w ten sam sposób, nawet jeżeli na pierwszy rzut oka nie wydawało się to możliwe. Nie chciałam także pisać czegoś, co zasadniczo byłoby trzykrotnie tą samą książką.

Aby identyfikować fragmenty istotne tylko dla danego języka\footnote{Ta konwencja pozostanie zachowana w miejscach, gdzie będzie mowa o wszystkich trzech językach w obrębie jednego rozdziału. Wszędzie indziej wszystkie zamieszczone informacje dotyczą języka polskiego.}, będą użyte kolory w następujący sposób:

\colorbox{cyan}{Česky}

\colorbox{yellow}{Po polsku}

\colorbox{magenta}{Po slovensky}

Wielu spośród członków Discorda Plovera, o których będę wspominać w Lekcji 1 jest bardzo dobrze obytych w komputerach. Niektórzy spośród nich budowali swoje własne komputery, znają się na programowaniu, i są najlepszym źródłem informacji dotyczących programu Plover. Na dobrą sprawę, wielu z nich używa stenografii właśnie do programowania.

Osobiście przez lata używałam oprogramowania Eclipse (również wspomniane w Lekcji 1) i aby być szczerą, wiem na jego temat znacznie więcej, niż na temat Plovera. Jeżeli nie jest się bardzo dobrze obeznanym w zakresie korzystania z komputera, ogromną pomocą mogłoby być znalezienie kogoś, kto pomoże zainstalować Plovera i dostosować go do używanego języka. Głównym językiem Plovera jest angielski, zwłaszcza w tematach technicznych.Po angielsku jest także dostępne wsparcie techniczne w programie Eclipse.

Niezależnie od tego, którego oprogramowania ostatecznie się używa, bardzo pomocna jest znajomość języka programowania Python.

Na niektórych stronach w książce znajdują się puste przestrzenie. Zostały pozostawione celowo, aby można było sporządzać swoje własne notatki. sleży zwrcać szczególną uwagę na adnotacje w ramkach, gdyż zawierają szczególnie istotne informacje.

%dedykacja
\clearpage
\begin{center}
\thispagestyle{empty}
\vspace*{\fill}
{\em Tę książkę dedykuję moim rodzicom.}
\vspace*{\fill}
\end{center}
\clearpage

\mainmatter

\chapter{Podstawy}

\pagebreak
\section{Lekcja 1 --- Czym jest stenografia maszynowa?}
\subsection{Ćwiczenie nr 1}

{\steno{\LARGE
\begin{tabular}{c c c c c c c c c c}
	X & S & T & V & L & -R & -L & G & -T & Y \\
	F & Z  & K  & P & R & C & B & -S & W & O \\
	J & L & C & E & R & -R & I & P & -L & V \\
	B & A & T & -S & K & G & U & Z & -T & S \\
	W & \texttildelow & X & O & F & Y & F & C & S & -L \\
	X & -R & Z & B & L & -L & R & -R & T & -T \\
	-S & S & E & B & S & W & Y & P & K & I \\
	\texttildelow & V & Z & G & -T & T & V & Y & C & P \\
	B & A & T & -S & K & G & U & Z & -T & S \\
	W & \texttildelow & X & O & F & Y & F & C & S & -L \\
	X & S & T & V & L & -R & -L & G & -T & Y \\
	F & Z  & K  & P & R & C & B & -S & W & O \\
	J & L & C & E & R & -R & I & P & -L & V \\
\end{tabular}}}

\newpage
\section{Lekcja 2 --- Zapoznanie z klawiaturą -- STREARST}

\subsection{Ćwiczenia}

Nie należy przejmować się tym, że poniższe kombinacje klawiszy nie tworzą prawdziwych słów.
Celem jest przyzwyczajenie Czytelnika do naciskania więcej niż jednego klawisza jednocześnie.
Należy pamiętać, że dywiz pomiędzy dwoma klawiszami wskazuje, że pierwszy należy nacisnąć lewą ręką, a drugi prawą.

\subsection{Ćwiczenie nr 2.1}

{\steno{\LARGE
	\begin{tabular}{c c c c c c c c}
		SE & AR & TE & AS & RE & AT & TA & ER \\
		RA & ES & SA & ET & S-R & EA & T-S & EA \\
		R-T & EA & AT & SA & AS & RA & AR & TA \\
		SE & ER & TE & ES & RE & ET & AS & SA \\
	\end{tabular}}}

\subsection{Ćwiczenie nr 2.2}

{\steno{\LARGE
	\begin{tabular}{c c c c c c}
		SET & TET & RET & SAT & TAT & RAT \\
		SER & TER & RER & SAR & TAR & RAR \\
		SES & TES & RES & SAS & TAS & SAT \\
	\end{tabular}}}

\newpage
\section{Lekcja 3 --- Zapoznanie z klawiaturą -- ZKVICLB}

\newpage
\section{Lekcja 4 --- Zapoznanie z klawiaturą -- XFPLJUGWOY} 

\newpage
\chapter{Części słów}
\pagebreak
\section{Lekcja 5 --- Akordy}
\newpage
\section{Lekcja 6 --- Klawisze a dźwięki}
\newpage
\section{Lekcja 7 --- Samogłoski}
\newpage
\section{Lekcja 8 --- Spółgłoski po lewej -- b, d, g, sz, cz, ż, dż}
\newpage
\section{Lekcja 9 --- Spółgłoski po lewej -- c, dz, ch, h, rz}
\newpage
\section{Lekcja 10 --- Spółgłoski po lewej -- m, n}
Niniejszą lekcję zawdzięczamy literom M i N.

\begin{tabela2}
	m & \steno{KP} \\
	n & \steno{LR} \\
\end{tabela2}

\section{Lekcja 11 --- Spółgłoski po lewej -- dź, ć, ś, ź, ń, ł}
\newpage
\section{Lekcja 12 --- Spółgłoski po prawej -- p, f, dz, j}
\newpage
\section{Lekcja 13 --- Spółgłoski po prawej -- d, z, k, ch, h}
\newpage
\section{Lekcja 14 --- Spółgłoski po prawej -- sz, cz, rz, ż, dż}
\newpage
\section{Lekcja 15 --- Spółgłoski po prawej -- m, n}
\newpage
\section{Lekcja 16 --- Spółgłoski po prawej -- dź, ć, ś, ź, ń, ł}
\newpage

\chapter{Kombinacje}
\pagebreak
\section{Lekcja 17 --- Kombinacje spółgłosek --- Kombinacje po lewej kończące się J oraz R}
\newpage
\section{Lekcja 18 --- Kombinacje po lewej kończące się l oraz ł}
\newpage
\section{Lekcja 19 --- R, rz, l oraz ł zgłoskotwórcze}
\newpage
\section{Lekcja 20 --- Kombinacje po lewej kończące się n, p, oraz b}
\newpage
\section{Lekcja 21 --- Kombinacje po lewej kończące się w, cz, k oraz m}
\newpage
\section{Lekcja 22 --- Kombinacje po lewej kończące się g, h, ch, d, dź, t, ć oraz sz}
\newpage
\section{Lekcja 23 --- Kombinacje po lewej kończące się z, ź, ż, s, c, dz oraz f}
\newpage
\section{Lekcja 24 --- Kombinacje po prawej --- Jedno i dwuklawiszowe kombinacje po prawej}
\newpage
\section{Lekcja 25 --- Trzyklawiszowe kombinacje po prawej}
\newpage
\section{Lekcja 26 --- Czteroklawiszowe kombinacje po prawej}
\newpage
\section{Lekcja 27 --- Trzygłoskowe kombinacje po lewej --- trzygłoskowe kombinacje po lewej kończące się j, l, ł, r, rz, n oraz ń}
\newpage
\section{Lekcja 28 --- Trzygłoskowe kombinacje po lewej kończące się b, t, ć, d, dź, k, g, ch, p, m, cz, dż, w, s oraz ź}

\newpage
\section{Lekcja 29 --- Czterogłoskowe kombinacje po prawej stronie}

\newpage
\section*{Gratulacje!}
Z pewnością można pogratulować sobie przebrnięcia przez tę sekcję książki. 
Jest to stanowczo najbardziej skomplikowany spośród rozdziałów, i mam nadzieję, że po jego lekturze
zrozumiałym jest, jak bardzo istotna jest umiejętność kodowania tak różnorodnych słów i ich części. Codzienna praktyka prawdopodobnie nie będzie wymagać używania wszystkich wymienionych kombinacji, jednakże można po nie sięgać w przyszłości w razie potrzeby. W stenografii zawsze lepiej jest być zbyt dobrze przygotowanym. Być może w tym momencie Czytelnik nie opanował jeszcze wszystkich kombinacji --- jest to zupełnie zrozumiałe; przyjęcie tak dużej ilości informacji wymaga czasu. W przyszłości przydatnym może być powtórzenie co bardziej złożonych konceptów w tej sekcji razem z przerabianiem dalszych partii materiału.

Teraz, posiadając dogłębne zrozumienie budowy części słów, z pewnością budowanie całych słów nie będzie stanowić dużego problemu. Następny rozdział jest znacznie prostszy niż ten --- powinna być to całkiem miła odmiana. W tym momencie można zrobić przerwę i odpocząć przed podjęciem dalszej nauki --- z pewnością się należy!
\newpage

\chapter{Przyrostki i przedrostki}
\section{Lekcja 30 --- Budowanie słów i przyrostki --- Przyrostki rzeczownikowe i przymiotnikowe}
\newpage
\section{Lekcja 31 --- Przyrostki czasownikowe}
\newpage
\section{Lekcja 32 --- Przyimki i przedrostki}
\newpage
\section{Lekcja 33 --- Trzy złośliwe przyimki niesylabiczne}
\newpage
\section{Lekcja 34 --- Wrostki}
\newpage
\chapter{Skrócenia}
\newpage
\section{Lekcja 35 --- Skrócenia, drugorzędne samogłoski, zaimki}
\newpage
\section{Lekcja 36 --- Zaimki pytające i wskazujące}
\newpage
\section{Lekcja 37 --- Liczby}
\newpage
\section{Lekcja 38 --- Skróty}
\newpage
\section{Lekcja 39 --- Alfabet i literowanie}
\newpage
\section{Lekcja 40 --- Zaimki złożone i samogłoski trzeciorzędne}
\newpage
\section{Lekcja 41 --- Redukcja akordów}
\newpage
\section{Lekcja 42 --- Zapożyczenia}
\newpage
\chapter{Słowo końcowe}
Ten tekst jest wyczerpującym opracowaniem na temat stenografi maszynowej w językach zachodniosłowiańskich. Mam nadzieję, ze Czytelnik nie spieszył się podczas czytania i dał sobie czas na przyswojenie treści. Być może okazało się, że rozum prędzej potrafi pojąć zasady niż palce --- a sytuacja, gdy dłonie i umysł nie są zsynchronizowane potrafi być bardzo frustrująca. Niech Czytelnik nie pozwoli, by tak było w jego przypadku. Najlepsze wyniki osiągnie się, przyswajając informacje z niniejszej książki w bardziej umiarkowanym tempie.

Sugeruję powracać do tej książki i regularnie powtarzać zawarte w niej ćwiczenia. Rekomenduję także ćwiczenia z pomocą tabel. Przechodzenie spokojne poprzez zawartość tabeli; pisanie wszystkiego z wyjątkiem wybranej kolumny; skoncentrowanie się na konkretnej kombinacji klawiszy; pisanie dwóch kombinacji spółgłosek na zmianę; przebiegnięcie przez wszystkie samogłoski z wszystkimi kombinacjami spółgłosek --- tego typu ćwiczenia pozwalają rozwinąć sprawność i wytrzymałość, co może być kluczowe w przypadku używania stenografii do, przykładowo, programowania czy długich sesji tłumaczeń.

Chciałabym powrócić do tego, jak bardzo ważne jest opanowanie pierwszych 34 lekcji w książce. Lekcje 35-42 są dodatkowym materiałem, który można włączyć po osiągnięciu pewnego poziomu sprawności w stenografowaniu. Uwaga --- informacje zawarte w rozdziale Skrócenia należy włączać małymi porcjami. Ponieważ te akordy stanowią wyjątki, przyswojenie ich zajmie więcej czasu. Jednakże jak tylko zostaną opanowane, ich użycie pozwoli znacznie zwiększyć prędkość pisania.

Sugeruję ćwiczyć tworząc dokładne kopie tektów pisanych celem zwiększenia dokładności. Pomoże to opanować palcowanie i rozkład języka na akordy bez ciśnienia wynikającego z obecności mówcy. Po opanowaniu tej umiejętności, sugeruję zacząć nagrywać teksty --- tempo nagrań powinno być dość spokojne. Nadążenie ze stenografowaniem za normalnym tempem ludzkiej mowy może okazać się szokująco trudne. Profesjonalni stenografowie zaczynają od tempa w okolicy 60 sylab na minutę, powoli przyśpieszając, gdy zyskują niezbędne po temu umiejętności. Dobrym pomysłem może być stenografowanie tekstów piosenek --- ballady i spokojny pop są wykonywane w tempie, które nie powinno być przytłaczające. Jednak należy się trzymać z dala od rapu --- może mieć bardzo, bardzo szybkie pasaże, z którymi nawet doświadczony stenograf będzie miał problem. Innym dobrym źródłem materiałów do ćwiczeń są książki dla dzieci. Można także poszukać tekstów dyktowanych w odpowiednim tempie na YouTube czy innych stronach internetowych.


Ta książka jest pierwszą z serii poświęconej stenografi maszynowej w językach słowiańskich. Kolejna będzie o językach południowosłowiańskich. Życzę powodzenia w przygodzie ze stenografią.

{\em Priscilla}
\appendix
\chapter{Dodatki}
\section{Dodatek A: Tabela liter}
\newpage
\section{Dodatek B: Przyrostki}
\newpage
\section{Dodatek C: Samogłoski -- po lewej, na środku, po prawej}
\newpage
\section{Dodatek D: Alfabetyczna lista grup spółgłoskowych po prawej stronie}
\newpage
\section{Dodatek E: Przyimki wewnętrzne}
\newpage
\section{Dodatek F: Złożone przyrostki}
\newpage
\section{Dodatek G: Skrócenia wewnętrzne}
\newpage
\section{Dodatek H: Interpukcja, akordy służące formatowaniu oraz identyfikacji mówiącego}

\end{document}
